\documentclass[10pt,a4paper]{amsart}
\usepackage[margin=19mm]{geometry}
\usepackage{amsmath,amssymb,mathtools}
\usepackage[scr=rsfs,cal=euler]{mathalfa}
\usepackage{xfrac}
\usepackage[unicode,hidelinks,bookmarks=false]{hyperref}
\newcommand{\ie}{i.e.\ }
\newcommand{\cf}{cf.\ }
\newcommand{\resp}{respectively\ }
\newcommand{\Subsection}{\subsection}
\providecommand{\nobreakdash}{\nobreak\hbox{-}}
\setlength{\emergencystretch}{2em}
\multlinegap0pt
\usepackage{amssymb}
\usepackage{mathtools}
\usepackage[shortlabels]{enumitem}
\newtheorem{lemma}{Lemma}
\newcommand{\G}{\mathcal G}
\title{New Extremal Ranges and Constructions of the Erd\H{o}s--Kleitman Problem}
\author{Cheng Chi, Yan Wang}
\date{Source: arXiv:2605.09535v3 (CC BY 4.0)}
\begin{document}
\maketitle

\noindent\emph{Source and licence.} New Extremal Ranges and Constructions of the Erd\H{o}s--Kleitman Problem. Version arXiv:2605.09535v3 (CC BY 4.0), distributed by arXiv under the Creative Commons Attribution 4.0 International licence, http://creativecommons.org/licenses/by/4.0/, as recorded in the arXiv licence metadata for this exact version and shown on the abstract page https://arxiv.org/abs/2605.09535v3. Full licence notice: \texttt{licenses/arxiv-2605.09535-CC-BY-4.0.txt}.\par\smallskip

\noindent\textit{Editorial scope.} Counted unit: the article's lemma lem:p-ordered-local-blocker (source0.tex lines 508--530). The article's complete original proof of it is delivered with it (source0.tex lines 533--686, one \texttt{proof} environment of 154 lines). No mathematics is written, repaired or re-derived: the statement and the proof are the article's own lines, in the article's own notation, and no retained line is substituted or re-flowed.  No further source-local context is needed: every cross-reference of every retained line resolves inside the two delivered spans.  Nothing was omitted from any retained span, so the package carries no omission record.  No retained line contains a citation, so the wrapper carries no bibliography and no bibliography entry.  No retained line contains a figure environment, an image-inclusion command or drawing code, so no image is delivered and the package declares no asset dependency.  Editorial formatting only, in addition: an AMS \texttt{amsart} preamble that restates the article's own declarations and macros used by the retained lines, namely the environments \texttt{lemma} on separate counters, together with the standard packages those lines need.  The retained environments print in this document as Lemma 1 in order of appearance, whereas the article prints the same items with the numbering of its own full document.  The complete native source of the article as distributed in the arXiv e-print for this exact version is bundled unchanged as \texttt{sources/200231/source0.tex}.
\begin{lemma}\label{lem:p-ordered-local-blocker}
        Let $m,p$ be integers with $m\ge2$ and $1\le p\le m-1$, and set
        $b=m+1-p$. There exist constants
        $\gamma_{\ref{lem:p-ordered-local-blocker}}>0$,
        $\rho_{\ref{lem:p-ordered-local-blocker}}>0$, and
        $C_{\ref{lem:p-ordered-local-blocker}}$ depending only on
        $m,p$ such that the following holds.

        Let $u,h$ be integers with
        $u\ge C_{\ref{lem:p-ordered-local-blocker}}$ and
        $0\le h\le\rho_{\ref{lem:p-ordered-local-blocker}}u$.
        Let $P,X$ and the families $\G_e^{(b)},\G_e^{(b+1)}$ be as in the definition above.
        If no mixed ordered partition exists, then
        \[
                \sum_{e\in\binom Pp}
                \left(
                \left|\binom Xb\setminus\G_e^{(b)}\right|
                +
                \left|\binom X{b+1}\setminus\G_e^{(b+1)}\right|
                \right)
                \ge \gamma_{\ref{lem:p-ordered-local-blocker}}(h+1)u^m.
        \]
\end{lemma}

\begin{proof}
        For each $e\in\binom Pp$, write $\mathcal B_e^{(b)}=\binom Xb\setminus\G_e^{(b)}$ and $\mathcal B_e^{(b+1)}=\binom X{b+1}\setminus\G_e^{(b+1)}$.
        Also write
        \[
                \mathcal Q_b=\{(e,T):e\in\binom Pp,\ T\in\mathcal B_e^{(b)}\} \text{ and }\mathcal Q_{b+1}=\{(e,T):e\in\binom Pp,\ T\in\mathcal B_e^{(b+1)}\}.
        \]
        We choose $\rho_{\ref{lem:p-ordered-local-blocker}}\le1/4$, and later choose $C_{\ref{lem:p-ordered-local-blocker}}$ sufficiently large.
        All constants below depend only on $m$ and $p$.

        First suppose $h=0$.
        Let $\mathcal P_0$ be the set of all ordered pairs of partitions $P=e_1\sqcup\cdots\sqcup e_u$ and $X=T_1\sqcup\cdots\sqcup T_u$, where $|e_r|=p$ and $|T_r|=b$ for every $r$.
        Since no mixed ordered partition exists, every member of $\mathcal P_0$ contains at least one forbidden pair from $\mathcal Q_b$.

        Let $I_0=\binom{pu}{p}\binom{bu}{b}$ be the number of possible $(p,b)$-pairs $(e,T)$.
        Fix such a pair $(e,T)$.
        Let $M_0$ be the number of members of $\mathcal P_0$ in which $e_r=e$ and $T_r=T$ for some $r\in[u]$.
        Then
        \[
                M_0
                =
                u\frac{(p(u-1))!}{(p!)^{u-1}}\frac{(b(u-1))!}{(b!)^{u-1}}
                =
                \frac{|\mathcal P_0|u}{I_0}.
        \]
        Indeed, one may first choose the common position $r$, then partitions $P\setminus e$ into $u-1$ ordered $p$-sets and $X\setminus T$ into $u-1$ ordered $b$-sets.

        Since every member of $\mathcal P_0$ contains at least one forbidden pair from $\mathcal Q_b$, we have $|\mathcal Q_b|M_0\ge |\mathcal P_0|$.
        Therefore,
        \[
                |\mathcal Q_b|
                \ge
                \frac{|\mathcal P_0|}{M_0}
                =
                \frac{I_0}{u}
                =
                \frac1u\binom{pu}{p}\binom{bu}{b}
                \ge c_0u^m
        \]
        for a constant $c_0=c_0(m,p)>0$, since $p+b=m+1$ and $u\ge C_{\ref{lem:p-ordered-local-blocker}}$.
        This proves the result when $h=0$.

        Now assume $h\ge1$.
        A \emph{base configuration} is a triple $\mathcal C_0=(\mathbf e,O,\mathbf T)$, where $P=e_1\sqcup\cdots\sqcup e_u$ with $|e_r|=p$, $O\in\binom Xh$, and $X\setminus O=T_1\sqcup\cdots\sqcup T_u$ with $|T_r|=b$.
        Let $\mathcal C$ be the set of all base configurations.

        For a base configuration $\mathcal C_0=(\mathbf e,O,\mathbf T)$, call $(e_r,T_r)$ the $r$-th position.
        Thus each position pairs the $r$-th $p$-block of $P$ with the $r$-th $b$-block of $X\setminus O$.
        We also write this position as $(r,e_r,T_r)$ to include its index.
        Here $e_r\subseteq P$ and $T_r\subseteq X\setminus O$.
        The overflow set $O$ is used only later, when a position is upgraded by replacing $T_r$ with $T_r\cup\{z\}$ for some $z\in O$.

        A position $(r,e_r,T_r)$ of $\mathcal C_0$ is called \emph{$b$-forbidden} if $T_r\in\mathcal B_{e_r}^{(b)}$.
        First suppose that at least $|\mathcal C|/2$ base configurations have at least $h+1$ $b$-forbidden positions.

        Let $I_b=\binom{pu}{p}\binom{bu+h}{b}$ be the number of possible $(p,b)$-pairs.
        Fix such a pair $(e,T)$. The number of base configurations in which $(e,T)$ occurs as a position $(e_r,T_r)$ is
        \[
                M_b
                =
                u\frac{(p(u-1))!}{(p!)^{u-1}}\binom{bu+h-b}{h}\frac{(b(u-1))!}{(b!)^{u-1}}
                =
                \frac{|\mathcal C|u}{I_b}.
        \]
        Indeed, one chooses the common position $r$, partitions $P\setminus e$, chooses the overflow set $O\subseteq X\setminus T$ of size $h$, and then partitions $X\setminus(O\cup T)$ into $u-1$ ordered $b$-sets.
        Hence, $|\mathcal Q_b|M_b\ge\frac12|\mathcal C|(h+1)$, and therefore
        \[
                |\mathcal Q_b|
                \ge
                \frac{|\mathcal C|(h+1)I_b}{2|\mathcal C|u}
                =
                \frac{h+1}{2}\frac{I_b}{u}
                \ge
                c_1(h+1)u^m,
        \]
        since $I_b\ge c'_1u^{m+1}$.
        This gives the desired conclusion in this case.

        We may therefore assume that more than half of the base configurations have at most $h$ $b$-forbidden positions.
        Call such configurations good.

        Fix a good base configuration $\mathcal C_0=(\mathbf e,O,\mathbf T)$.
        Define the bipartite overflow graph $\Gamma_O(\mathcal C_0)$ as follows.
        One part consists of the positions $(r,e_r,T_r)$, and the other part is $O$.
        Join a position $(r,e_r,T_r)$ to $z\in O$ if $T_r\cup\{z\}\in\G_{e_r}^{(b+1)}$.
        Here an upgrade means replacing a $b$-set $T_r$ by $T_r\cup\{z\}$, where $z\in O$, so that the tail has size $b+1$.
        The upgrade is allowed precisely when $T_r\cup\{z\}\in\G_{e_r}^{(b+1)}$.
        Thus a missing edge is exactly a forbidden upgrade.

        We claim that $\Gamma_O(\mathcal C_0)$ has at least $h$ missing edges.
        Suppose to the contrary that it has less than $h$ missing edges.
        We first show that the $b$-forbidden positions can be assigned distinct vertices of $O$ by allowed upgrades.
        In fact, if no such assignment exists, then Hall's condition would give a nonempty family $Y$ of $b$-forbidden positions with $|N(Y)|<|Y|$, where $N(Y)$ denotes the neighborhood of $Y$ in the overflow graph.
        Then all pairs in $Y\times (O\setminus N(Y))$ are missing edges, so the number of missing edges is at least
        \[
                |Y|(h-|N(Y)|)
                \ge |Y|(h-|Y|+1)
                \ge h,
        \]
        because $1\le |Y|\le h$.
        This contradicts the assumption that there are fewer than $h$ missing edges.

        Hence all $b$-forbidden positions can be matched injectively into $O$ by allowed upgrades.
        After doing this, assign the remaining overflow vertices greedily to unused positions.
        When a remaining $z\in O$ is considered, fewer than $h$ positions have already been used, and $z$ is non-adjacent to fewer than $h$ positions.
        Hence fewer than $2h$ positions are unavailable.
        Since $h\le u/4$, an unused neighbor remains.

        After all overflow vertices are assigned, upgrade exactly the positions to which they were assigned.
        Every upgraded position is allowed as a $(b+1)$-tail, and every not upgraded position is not $b$-forbidden, hence is allowed as a $b$-tail.
        This gives a mixed ordered partition, contradicting the hypothesis.
        Therefore every good base configuration has at least $h$ forbidden upgrades.

        We now count these forbidden upgrades.
        Let $I_{b+1}=\binom{pu}{p}\binom{bu+h}{b+1}$ be the number of possible $(b+1)$-pairs.
        Each base configuration has exactly $uh$ upgrade incidences.
        Fix a possible $(p,b+1)$-pair $(e,U)$. The number of upgrade incidences $(\mathcal C_0,r,z)$ realizing $(e,U)$, meaning $e_r=e$ and $T_r\cup\{z\}=U$, is
        \[
                M_{b+1}
                =
                u(b+1)\frac{(p(u-1))!}{(p!)^{u-1}}\binom{bu+h-b-1}{h-1}\frac{(b(u-1))!}{(b!)^{u-1}}
                =
                \frac{|\mathcal C|uh}{I_{b+1}}.
        \]
        Indeed, one chooses the common position $r$, chooses the overflow vertex $z\in U$, sets $T_r=U\setminus\{z\}$, partitions $P\setminus e$, chooses the remaining $h-1$ overflow vertices from $X\setminus U$, and partitions the remaining vertices of $X$ into $u-1$ ordered $b$-sets.
        Since more than half of the base configurations are good and every good configuration has at least $h$ forbidden upgrades, we have
        \[
                |\mathcal Q_{b+1}|
                M_{b+1}
                \ge
                \frac12|\mathcal C|h,
        \]
        and hence
        \[
                |\mathcal Q_{b+1}|
                \ge
                \frac{I_{b+1}}{2u}
                \ge c_2 u^{m+1}.
        \]
        Since $h\le \rho_{\ref{lem:p-ordered-local-blocker}}u\le u/4$, we have $h+1\le u$, and hence $|\mathcal Q_{b+1}| \ge c_2(h+1)u^m$.
        Taking $\gamma_{\ref{lem:p-ordered-local-blocker}}\le \min\{c_0,c_1,c_2\}$ proves
        \[
                |\mathcal Q_b|+|\mathcal Q_{b+1}|
                =
                \sum_{e\in\binom Pp}
                \left(
                \left|\binom Xb\setminus\G_e^{(b)}\right|
                +
                \left|\binom X{b+1}\setminus\G_e^{(b+1)}\right|
                \right)
                \ge
                \gamma_{\ref{lem:p-ordered-local-blocker}}(h+1)u^m.
        \]
        This completes the proof.
\end{proof}

\end{document}
